Het was weer een maandag.
Maandag, de dag van introspectie en reflectie, past bij de teruggetrokken aard van Sophia.

Geen van haar relaties
kende haar emoties.

Ze zag die als onecht.
Deed het een ander recht

als een paard steigerde
en voortgaan weigerde

de grillen van het dier
te zien in de koetsier?

Slechts de onbekwaamheid
om de ontembaarheid

mild te beteugelen,
liet zonden vleugelen.

Wegens dit grote hart
noemde men haar apart

of gewoon Sophia.
“Soof” op privépodia.

Sophia is wijsheid

Ja, Sophia betekent inderdaad wijsheid. Het is een Grieks woord (σοφία) dat in veel filosofische, theologische en mystieke tradities een diepe betekenis heeft. In het bijzonder:

In de Griekse filosofie staat sophia voor goddelijke of hogere wijsheid, in tegenstelling tot phronesis (praktische wijsheid).

In het christendom wordt Sophia soms gepersonifieerd als de goddelijke wijsheid zelf, zoals in de Wijsheidsboeken van het Oude Testament (bijv. Wijsheid van Salomo).

In gnosticisme is Sophia een centrale figuur: een vrouwelijke emanatie van God die door haar val de materiële wereld mogelijk maakte — maar ook de weg terug naar het goddelijke toont.

In oosterse orthodoxie wordt Sophia soms geassocieerd met de Heilige Geest of met Christus als Wijsheid van God.

In jouw context, waarin Sophia de personificatie is van spirituele luiheid die zich voordoet als doofstom en hard werkt als naaister, roept dat een boeiende paradox op: de naam suggereert wijsheid, maar de rol toont iemand die zich onttrekt aan innerlijk weten door uiterlijke arbeid en schijnbare verstilling. Dat kan gelezen worden als een ironische of symbolische laag — wijsheid die sluimert, verdraaid wordt of zich pas toont wanneer ze haar escapisme opgeeft.

Wil je dat haar naam uiteindelijk waar wordt? Dat ze dus ook daadwerkelijk Sophia wordt — wijsheid die zich uit de slaap ontwaakt?

Schaarse acceptatie,
gebrek aan empathie,

valse dominatie
en leerimitatie

waren de stadswerken.
Deze stijlkenmerken

kweekten bij Sophia
vlot de paranoia,

geen aandacht te trekken,
geen haat op te wekken,

geen scepsis te spekken,
geen kans te verstrekken

om haar te judassen.
Stil ging ze verkassen

als ze wat onraad rook.
Stil en doof was ze ook

in letterlijke zin.
Toch niet in het begin

toen ze dicht bij God stond.
Haar hart was niet verwond

toen ze op aarde kwam,
zonder woord, zonder schram,

maar met de zegening
te doen een voorspelling

over ieders toekomst.
De stad bestreed aartstomst

de verleende gave
zonder hem te staven.

Ze stichtte verwarring
met elke voorspelling.

Bovenal het onheil,
dat door ieders leefstijl

boven de stadstaat hing,
was een loze melding.

Ze heette Sophia,
maar deed als Cassandra,

die ook niet werd geloofd.
Ze werd dito onthoofd

in symbolische zin.
Ze hield haar woorden in.

Men kon haar niet laken
of verwijten maken

over een opmerking
zonder een taaluiting.

Niemand vond dat verdacht.
Niemand gaf haar aandacht.

Maar dat bracht ook mede
een veilige vrede.

Geen verwijt en geen smart.
Uit het oor, in het hart.

Ze ontving meer fatsoen
door gewoon goed te doen.

Geen woorden maar daden
temperde de kwaden.

Zoals de broodbakker.
Zijn zoon was een stakker.

Vrij zwaar gehandicapt.
Waar zij niet over rept,

maar wel een oog voor had.
Zij nam hem mee op pad

in een houten rolstoel
voorbij de kikkerpoel.

Ze vermaakte kleding
wegens zijn beperking.

De geharde dikzak
kreeg aldoende een zwak

voor onze Sophia,
zijn moeder Maria.

Als je toch iets wil doen ga dan niets doen

Opzij, opzij opzij van Herman Veen

Met druk, druk, druk, zo lui.
vermeed men elke bui.
De kunst van zitten en nietsdoen

Je wordt niets van waarde door het wel of niet resultaat maar door het lef ervoor te gaan
